% ============================================================================
% sec3_structural_model.tex — REVISION-V6: Full structural model (RAND reframe)
% Paper: Frequent Losers as Cover Bidders in Public Procurement
% Target: RAND Journal of Economics
% ============================================================================

\section{A Framework for Cover-Bidder Deployment} % REVISION-V6: M3 — renamed from "Structural Model"
\label{sec:structural}

We develop a framework for cover bidding in procurement auctions
that formalizes the cover-bidder deployment decision, characterizes
equilibrium cover-bid distributions, and generates testable predictions
that map directly to our empirical strategy
(Section~\ref{sec:empirical}). The framework has three aims: (i)~to
provide economic foundations for the frequent-loser screen; (ii)~to
specify a mixture likelihood that separates cover bids from genuine bids
in the data; and (iii)~to derive the market-selection prediction---that
cover bidding concentrates in competitive markets---as an endogenous
outcome of the cartel's optimization.

% ── 3.1 Environment ──────────────────────────────────────────────

\subsection{Environment}
\label{sec:environment}

Consider a procurement auction for item $i$ at purchasing unit $k$ in
year $t$. The procuring agency posts a reference price $r > 0$ and
selects the lowest bidder. There are $n \geq 1$ \emph{genuine} bidders,
each with private cost $c_j \sim F(\cdot \mid \mathbf{x}_j)$, where
$\mathbf{x}_j$ is a vector of observable cost shifters (firm size, CNAE
sector, geographic distance to PBU~$k$). Genuine bidders submit bids
$b_j = c_j + \mu_j$, where $\mu_j \geq 0$ is a markup that depends on
the competitive environment.

A cartel controls firm $j^*$---the designated winner---and deploys
$m \geq 0$ \emph{cover bidders} (frequent losers). Each cover bidder
submits a bid $b_\ell$ drawn from a distribution $G(\cdot \mid b^*,
\boldsymbol{\theta})$, where $b^*$ is the cartel's designated winning
bid and $\boldsymbol{\theta}$ is a vector of strategic parameters. Cover
bids are designed to \emph{not} win: $b_\ell > b^*$ with probability
one. The procuring agency observes $n + m$ total bids.

Two institutional features of S\~ao Paulo's BEC platform shape the
model:

\begin{enumerate}
  \item \textbf{Convite} (sealed-bid): requires at least three bidders
    (Lei~8.666/93, Art.~22, \S7). If $n < 3$, the cartel must deploy
    $m \geq 3 - n$ cover bidders to satisfy the minimum-bidder
    constraint. Bids are submitted simultaneously and opened together.
  \item \textbf{Preg\~ao} (electronic reverse auction): no
    minimum-bidder requirement. Bidders submit descending bids in real
    time. Cover bidders can observe the auction dynamics before
    calibrating their bids, facilitating coordination.
\end{enumerate}

\noindent We denote the auction modality by $\tau \in \{\text{convite},
\text{preg\~ao}\}$, the minimum-bidder constraint by $\underline{n}(\tau)$
(equal to 3 for convite, 0 for preg\~ao), and the detection probability
by $\theta_k \in [0,1]$, which varies across PBUs and is decreasing in
oversight capacity.

% ── 3.2 Cartel's Problem ─────────────────────────────────────────

\subsection{The Cartel's Optimization Problem}
\label{sec:cartel_problem}

The cartel chooses the number of cover bidders $m$, the winning bid
$b^*$, and the cover-bid distribution $G$ to maximize expected profit
net of cover-bidding costs and detection penalties:
\begin{equation}
  \max_{m, b^*, G} \; \underbrace{(b^* - c_{j^*})}_{\text{cartel surplus}}
  - \underbrace{C(m)}_{\text{cover cost}}
  - \underbrace{\theta_k \cdot \Phi(m, b^*)}_{\text{expected penalty}}
  \label{eq:cartel_objective}
\end{equation}
subject to:
\begin{align}
  b^* &\leq r \label{eq:ref_price_constraint} \\
  \Pr(b_\ell > b^* \mid G) &= 1 \quad \forall \, \ell \in \{1, \ldots, m\}
    \label{eq:cover_above} \\
  n + m &\geq \underline{n}(\tau) \label{eq:min_bidder}
\end{align}

\noindent where $C(m) = c_0 + c_1 m$ is the linear cost of deploying $m$
cover bidders (recruitment, coordination, and compensation), and
$\Phi(m, b^*)$ is the penalty conditional on detection, which we assume
is increasing in $m$ (more cover bidders leave a larger forensic
footprint) and in the gap $b^* - c_{j^*}$ (a larger markup attracts
more scrutiny).

\begin{proposition}[Optimal number of cover bidders]
\label{prop:optimal_m}
Let $\pi(b^*, n) = b^* - c_{j^*}$ be the cartel's gross surplus. Under
the cost structure $C(m) = c_0 + c_1 m$ and detection penalty
$\Phi(m) = \phi_0 m$ (linear in $m$), the cartel's optimal cover-bidder
count is:
\begin{equation}
  m^* = \max\left\{\underline{n}(\tau) - n, \;
  \left\lfloor \frac{\partial \pi / \partial m}{c_1 + \theta_k \phi_0}
  \right\rfloor \right\}
  \label{eq:optimal_m}
\end{equation}
\noindent The first argument is the mechanical minimum to satisfy the
bidder-count constraint (binding only under convite when $n < 3$). The
second is the profit-maximizing number of cover bidders when the
constraint does not bind. The optimal $m^*$ is:
\begin{itemize}
  \item Increasing in $n$ if cover bidders help sustain the winning bid
    against genuine competition (i.e., $\partial^2 \pi / \partial m
    \partial n > 0$);
  \item Decreasing in $\theta_k$ (higher detection probability raises
    the marginal cost of cover bidding);
  \item Decreasing in $c_1$ (higher per-cover-bidder cost);
  \item Weakly larger under convite than preg\~ao, holding $n$ fixed,
    because the minimum-bidder constraint may bind.
\end{itemize}
\end{proposition}

\paragraph{Argument.}
The first-order condition of~\eqref{eq:cartel_objective} with respect
to $m$, ignoring the integer constraint, is $\partial \pi / \partial m
= c_1 + \theta_k \phi_0$. Since $\partial^2 \pi / \partial m^2 < 0$
(decreasing marginal returns to cover bidding), the interior solution is
unique. The binding constraint $n + m \geq \underline{n}(\tau)$
produces the corner solution $m = \underline{n}(\tau) - n$ when this
exceeds the interior solution. The comparative statics follow from the
implicit function theorem applied to the FOC.



% ── 3.3 Cover-Bid Distributions ──────────────────────────────────

\subsection{Equilibrium Cover-Bid Distributions}
\label{sec:cover_distributions}

The cartel chooses the cover-bid distribution $G$ to minimize the
probability of detection while satisfying constraint~\eqref{eq:cover_above}.
We characterize two equilibrium regimes that differ in the cartel's
information and coordination capacity.

\begin{proposition}[Regime~1: Complementary cover bidding]
\label{prop:regime1}
When the cartel has limited coordination capacity (it can instruct cover
bidders to ``bid above $b^*$'' but cannot specify exact bid levels), the
equilibrium cover-bid distribution is:
\begin{equation}
  G_1(b) = \frac{b - b^*}{\delta}, \quad b \in [b^*, b^* + \delta]
  \label{eq:regime1}
\end{equation}
where $\delta > 0$ is the maximum spread above the winning bid. Under
Regime~1, cover bids are uniformly distributed above $b^*$, producing
high bid dispersion ($\sigma_{\text{cover}} > \sigma_{\text{genuine}}$)
and a right-skewed bid distribution.
\end{proposition}

\paragraph{Argument.}
The detection-minimizing distribution over $[b^*, b^* + \delta]$
subject to $b > b^*$ and finite support is uniform, because any
concentration of mass (e.g., at a focal point above $b^*$) would be
detectable by a variance screen. The uniform distribution maximizes
entropy over the feasible support, minimizing the statistical
distinguishability of cover bids from genuine high bids.



\begin{proposition}[Regime~2: Coordinated cover bidding]
\label{prop:regime2}
When the cartel has full coordination capacity (it knows the winning bid
$b^*$ and can specify cover bids precisely), the equilibrium cover-bid
distribution is:
\begin{equation}
  G_2(b) = \Phi\!\left(\frac{b - \mu_c}{\sigma_c}\right),
  \quad \mu_c = b^* + \epsilon, \quad
  \sigma_c < \sigma_{\text{genuine}}
  \label{eq:regime2}
\end{equation}
where $\Phi(\cdot)$ is the standard normal CDF, $\epsilon > 0$ is a
small margin ensuring $b_\ell > b^*$, and $\sigma_c$ is chosen to mimic
the genuine bid distribution. Under Regime~2, cover bids are tightly
clustered near $b^*$, producing low bid dispersion
($\sigma_{\text{cover}} < \sigma_{\text{genuine}}$) and a bid
distribution that is harder to distinguish from genuine competition.
\end{proposition}

\paragraph{Argument.}
When the cartel observes $b^*$ (feasible under preg\~ao's real-time bid
revelation), it can instruct cover bidders to bid at $b^* + \epsilon$
with small noise. The optimal noise variance $\sigma_c^2$ trades off
detection risk (a degenerate distribution at $b^* + \epsilon$ is
instantly detectable) against the constraint that cover bids must not be
\emph{too} dispersed (which would reveal their non-genuine origin).
Under the assumption that the detection algorithm uses a bid-dispersion
screen (as in \citealt{imhof2019detecting}), the optimal
$\sigma_c^2$ matches $\sigma_{\text{genuine}}^2$ from below, producing
$\sigma_c < \sigma_{\text{genuine}}$.



% ── 3.4 Market Selection ─────────────────────────────────────────

\subsection{Market Selection: Why Cover Bidding Concentrates in
Competitive Markets}
\label{sec:market_selection}

\begin{proposition}[Market-selection prediction]
\label{prop:market_selection}
The optimal number of cover bidders $m^*$ is decreasing in the
Herfindahl--Hirschman Index (HHI) of genuine competition. Cover bidding
concentrates in competitive markets (low HHI) rather than concentrated
markets (high HHI).
\end{proposition}

\paragraph{Argument.}
In concentrated markets with high HHI, the cartel's designated winner
$j^*$ already faces few genuine competitors. The marginal value of an
additional cover bidder is low: with only one or two genuine rivals, the
cartel can sustain high prices through tacit coordination or market
power without needing to manufacture the appearance of competition. The
return $\partial \pi / \partial m$ is therefore small, and the FOC
$\partial \pi / \partial m = c_1 + \theta_k \phi_0$ implies $m^*
\approx 0$.

In competitive markets with low HHI, genuine bidders exert downward
pressure on the winning bid. Cover bidders serve two functions: (i)
creating noise that obscures the cartel's bid pattern from forensic
screens, and (ii) under convite, ensuring the minimum-bidder constraint
is satisfied despite genuine competitors potentially declining to bid.
The return $\partial \pi / \partial m$ is large, and the FOC implies
$m^* > 0$.

Formally, the cross-derivative $\partial^2 \pi / \partial m \partial
\text{HHI} < 0$: the marginal return to cover bidding decreases with
market concentration. By the implicit function theorem,
$\partial m^* / \partial \text{HHI} < 0$.



This proposition provides the theoretical grounding for the network-split
result (Section~\ref{sec:results_network}): the entire FL--price effect
concentrates among FL firms operating in competitive markets (low winner
HHI), while FL firms in concentrated markets show no significant price
effect. The model predicts this pattern as an equilibrium outcome, not a
data anomaly.

% ── 3.5 Structural Likelihood ────────────────────────────────────

\subsection{Structural Likelihood}
\label{sec:likelihood}

We specify a mixture model for the bid-level data. Each bid $b_{jt}$ in
tender $t$ is generated by one of two processes:

\begin{equation}
  b_{jt} \sim
  \begin{cases}
    F_{\text{genuine}}(b \mid \mathbf{x}_j, \boldsymbol{\alpha}) &
      \text{if firm } j \text{ is genuine (non-FL)} \\
    G_{\text{cover}}(b \mid b^*_t, \boldsymbol{\gamma}) &
      \text{if firm } j \text{ is FL}
  \end{cases}
  \label{eq:mixture}
\end{equation}

\noindent where $\boldsymbol{\alpha} = (\mu_g, \sigma_g^2)$ are the
parameters of the genuine bid distribution and $\boldsymbol{\gamma} =
(\delta \text{ or } \mu_c, \sigma_c)$ are the parameters of the
cover-bid distribution.

\paragraph{Genuine bids.} We model genuine bids as log-normal:
\begin{equation}
  \log b_{jt} \mid j \text{ genuine} \sim
  N\!\left(\mathbf{x}_j'\boldsymbol{\beta} + \alpha_i + \lambda_t,
  \; \sigma_g^2\right)
  \label{eq:genuine_dist}
\end{equation}
where $\alpha_i$ and $\lambda_t$ are item and year fixed effects, and
$\mathbf{x}_j$ includes firm size (porte), firm age, and CNAE sector
dummies.

\paragraph{Cover bids.} Under Regime~1:
\begin{equation}
  b_{jt} \mid j \text{ FL} \sim
  \text{Uniform}[b^*_t, \; b^*_t + \delta]
  \label{eq:cover_dist_r1}
\end{equation}
Under Regime~2:
\begin{equation}
  \log b_{jt} \mid j \text{ FL} \sim
  N\!\left(\log b^*_t + \epsilon, \; \sigma_c^2\right),
  \quad \sigma_c < \sigma_g
  \label{eq:cover_dist_r2}
\end{equation}

\paragraph{Likelihood.} The log-likelihood contribution of tender $t$
with $n_t$ genuine bids and $m_t$ FL bids is:
\begin{equation}
  \ell_t = \sum_{j \in \mathcal{G}_t} \log f_{\text{genuine}}(b_{jt}
  \mid \mathbf{x}_j, \boldsymbol{\alpha})
  + \sum_{\ell \in \mathcal{F}_t} \log g_{\text{cover}}(b_{\ell t}
  \mid b^*_t, \boldsymbol{\gamma})
  \label{eq:loglik}
\end{equation}
where $\mathcal{G}_t$ and $\mathcal{F}_t$ are the sets of genuine and
FL bidders in tender $t$, respectively. The FL classification provides
the mixture labels, so no EM algorithm is required---the model is a
supervised mixture with known group membership.

\paragraph{Estimation.} We estimate $(\boldsymbol{\alpha}, % R1-FIX: sample sizes corrected
\boldsymbol{\gamma})$ by maximum likelihood using the bid-level data
(approximately 40 million bids total; 23.2 million non-FL losing bids
and 189,381 FL losing bids after restricting to observations with valid
prices and item-year fixed-effect coverage). Standard errors are
computed via parametric bootstrap (500 replications, parallelized
across 16 CPU cores). The genuine-bid parameters
$(\boldsymbol{\beta}, \sigma_g^2)$ are estimated from the non-FL
losing bids; the cover-bid parameters
$(\delta \text{ or } \sigma_c)$ are estimated from the FL losing
bids, conditional on the winning bid $b^*_t$ in each tender.

\paragraph{Model selection.} We select between Regime~1 and Regime~2
using the Bayesian Information Criterion (BIC) and a Kolmogorov--Smirnov
goodness-of-fit test comparing the fitted cover-bid distribution to the
empirical FL bid distribution. The regime with lower BIC and smaller KS
distance is selected as the preferred specification.

% ── 3.6 Identification ───────────────────────────────────────────

\subsection{Identification of Structural Parameters}
\label{sec:identification}

Each structural parameter is identified by specific variation in the
data:

\begin{itemize}
  \item $\boldsymbol{\beta}$, $\sigma_g^2$ (genuine bid distribution):
    identified from 37.5 million non-FL bids within item--year cells.
    The item fixed effects absorb product-specific cost levels; the year
    fixed effects absorb aggregate trends. Firm-level covariates
    $\mathbf{x}_j$ provide cross-sectional variation in costs.

  \item $\delta$ or $\sigma_c$ (cover bid parameters): identified from
    the distribution of FL bids \emph{relative to the winning bid} in
    each tender. The model requires that FL bids are measured jointly
    with $b^*_t$, which we observe directly as the negotiated price.

  \item $\theta_k$ (detection probability): partially identified from
    the cross-section of PBU characteristics. Following
    Proposition~\ref{prop:optimal_m}, $m^*$ is decreasing in
    $\theta_k$. We proxy $\theta_k$ with PBU size (number of
    procurement staff), exploiting the oversight heterogeneity
    documented in Section~\ref{sec:robustness}.

  \item The \emph{implied markup}
    $\hat{\mu} = E[b^* - c_{j^*}] / E[c_{j^*}]$ is a structural
    quantity that can be compared with the reduced-form OLS coefficient
    (6.4\%) and the IV estimate (21\%) as a consistency check.
\end{itemize}

% ── 3.7 Testable Predictions ─────────────────────────────────────

\subsection{Testable Predictions}
\label{sec:predictions}

The model generates five testable predictions, each mapped to a specific
empirical test in Section~\ref{sec:empirical}:

\begin{prediction}[Price effect]
\label{pred:price_v6}
FL-present tenders have higher winning prices than comparable tenders
without FL, controlling for item type, year, and PBU.
\emph{Test:} OLS regression (Section~\ref{sec:reduced_form}).
\end{prediction}

\begin{prediction}[Strategic selection]
\label{pred:selection_v6}
FL-present tenders have at least as many genuine (non-FL) competitors as
FL-absent tenders. Cover bidders are deployed on top of genuine
competition, not as substitutes.
\emph{Test:} Non-FL firm count regression (Section~\ref{sec:reduced_form}).
\end{prediction}

\begin{prediction}[Regime identification]
\label{pred:regime_v6}
Under Regime~1, FL bid dispersion exceeds genuine bid dispersion
($\sigma_{\text{cover}} > \sigma_{\text{genuine}}$). Under Regime~2,
FL bid dispersion is lower ($\sigma_{\text{cover}} <
\sigma_{\text{genuine}}$). BIC model selection determines the
empirically dominant regime.
\emph{Test:} Structural likelihood comparison and bid-dispersion
regression (Section~\ref{sec:structural_estimation}).
\end{prediction}

\begin{prediction}[Market selection]
\label{pred:market_v6}
The FL--price effect is decreasing in market concentration (HHI).
Cover bidding concentrates in competitive markets where the marginal
return to manufactured competition is highest.
\emph{Test:} Network-split regression and FL $\times$ HHI interaction
(Section~\ref{sec:reduced_form}).
\end{prediction}

\begin{prediction}[Exchangeability violation]
\label{pred:exchange_v6}
FL bid residuals are non-exchangeable with non-FL bid residuals
conditional on observables. The structural model implies that FL bids
are drawn from $G_{\text{cover}}$ rather than $F_{\text{genuine}}$,
producing distinguishable residual distributions.
\emph{Test:} Kolmogorov--Smirnov test on Bajari--Ye first-stage
residuals (Section~\ref{sec:bajari_ye}).
\end{prediction}

% REVISION-V6: Prediction 6 (minimum-bidder RDD) removed — BEC records
% unit prices, not total contract values, so the R$80K modality threshold
% cannot serve as an RDD running variable with available data.
